\section{Адаптация к условиям экономики России}
Российская спецификация строится как концептуальная адаптация модели Sims--Zha, а не как точное воспроизведение американских рядов. Базовая выборка начинается с 2014 года, поскольку в конце 2013 --- начале 2014 года Банк России перешёл к режиму инфляционного таргетирования, в рамках которого ключевая ставка стала центральной переменной денежно-кредитной политики. Поэтому более ранний период не включается в базовую оценку, так как он относится к иному режиму ДКП.

Итоговый вектор переменных в оценённой российской спецификации задаётся как

$$
X_t =
\left[
\log P_{cm,t},
\log M2_t,
R_t,
\log P_{im,t},
\log P_t,
\log W_t,
\log Y_t,
\log Debt_t
\right]'.
$$

Такой порядок не повторяет запись переменных в тексте статьи один к одному, но соответствует идентифицирующей логике оценённой SVAR. Сначала ставится быстро реагирующий индикатор сырьевых цен $P_{cm,t}$, затем денежный блок $(M2_t, R_t)$, после чего идут более инерционные переменные частного сектора.

В российской спецификации используется одна модель с ключевой ставкой и денежным агрегатом M2. Такой выбор связан с тем, что после перехода Банка России к инфляционному таргетированию ключевая ставка стала центральным инструментом ДКП, а её решения и прогнозы влияют на ожидания участников рынка; поэтому спецификация с резервами банковской системы не используется \cite{tishin2019monetary,abdurakhmanov2023keyrate}.

В спецификации используется VAR(1). Это сознательная упрощённая спецификация: информационные критерии на короткой российской выборке склоняют к большему числу лагов, но базовая оценка с одним лагом выбрана как более компактная и устойчивая для структурной идентификации и интерпретации импульсных откликов.

\paragraph{Источники данных.}
Для выпуска $Y_t$ используется квартальный сезонно скорректированный ряд реального ВВП России из базы Института экономической политики имени Е.,Т. Гайдара \cite{iep_gdp_sa}. Для уровня цен $P_t$ используется индекс потребительских цен ЕМИСС; цепной индекс переводится в ряд уровней \cite{fedstat_cpi_31074}. Заработная плата $W_t$ задаётся через среднюю месячную заработную плату Росстата, приведённую к реальному выражению с помощью ИПЦ и агрегированную до квартальной частоты \cite{rosstat_short_term_indicators}.

Денежно-кредитный блок включает ключевую ставку Банка России $R_t$ и денежный агрегат M2. Ключевая ставка берётся из статистики Банка России и используется в процентных пунктах без логарифмирования \cite{cbr_key_rate}. Для денежной массы используется сезонно скорректированный ряд M2 Банка России \cite{cbr_m2_sa}.

Переменная $Debt_t$ строится по данным Банка России о просроченной задолженности по кредитам юридических лиц-резидентов и индивидуальных предпринимателей. Для периода после 2019 года используется сводная таблица Банка России, для более раннего периода --- отдельные таблицы по рублёвой задолженности и задолженности в иностранной валюте и драгоценных металлах \cite{cbr_overdue_debt_after_2019,cbr_overdue_debt_rub_before_2019,cbr_overdue_debt_fx_before_2019}. Этот показатель используется как российский аналог переменной банкротств в исходной модели.

Индекс $P_{cm}$ строится как сырьевой ценовой индекс на основе цен производителей в добыче полезных ископаемых \cite{fedstat_ppi_43561,fedstat_ppi_57609}. Такая трактовка сохраняет логику исходной модели Sims--Zha, где сырьевые цены используются как быстро реагирующий индикатор инфляционного давления.

Индекс $P_{im}$ отражает цены промежуточных материалов и строится по узкому набору промежуточных отраслей. Выбор отраслей опирается на группировку, использованную Добронравовой для российских отраслей промышленности: с опорой на классификации Ganley и Salmon, а также Hayo и Uhlenbrock, она относит к кластеру промежуточных товаров $PRIM$ производство кокса и нефтепродуктов, химическое производство, металлургию и производство готовых металлических изделий \cite{dobronravova2022industry}. В российской спецификации этим отраслям соответствуют коды $DF$, $DG$ и $DJ$ в ОКВЭД и коды 19, 20, 24 и 25 в ОКВЭД2. Для построения индекса используются отраслевые индексы цен производителей ЕМИСС и веса Росстата для расчёта ИЦП по видам экономической деятельности \cite{fedstat_ppi_43561,fedstat_ppi_57609,rosstat_ppi_weights}.

Все ряды приводятся к квартальной частоте. Все ряды, кроме ключевой ставки, приведены в сезонно скорректированном виде. Все уровневые переменные логарифмируются, кроме ключевой ставки, которая сохраняется в процентных пунктах.
